\chapter{KẾT LUẬN VÀ HƯỚNG PHÁT TRIỂN}

\section{Kết quả đạt được}
Trong quá trình thực hiện đề tài \textbf{"Nghiên cứu và xây dựng hệ thống KMA Chatbot hỗ trợ tư vấn tuyển sinh"}, nhóm nghiên cứu đã nỗ lực tìm hiểu, phân tích và áp dụng các công nghệ tiên tiến trong lĩnh vực Xử lý ngôn ngữ tự nhiên (NLP) để giải quyết bài toán tư vấn tuyển sinh tại Học viện Kỹ thuật Mật mã. Những kết quả nổi bật mà đề tài đã đạt được bao gồm:

\begin{itemize}
    \item \textbf{Về mặt nghiên cứu lý thuyết:} Đã tìm hiểu sâu sắc về kiến trúc hệ thống Chatbot dựa trên nền tảng mã nguồn mở Rasa, bao gồm Rasa NLU và Rasa Core. Nắm bắt được phương pháp kết hợp mô hình ngôn ngữ tiếng Việt (PhoBERT) để nâng cao độ chính xác trong việc nhận diện ý định (Intent) và trích xuất thực thể (Entity). Đồng thời, đề tài đã nghiên cứu thành công kỹ thuật RAG (Retrieval-Augmented Generation) để giải quyết giới hạn của chatbot truyền thống khi cần tra cứu các tài liệu dài và chuyên sâu.
    
    \item \textbf{Về mặt phát triển hệ thống:} Đã xây dựng thành công kiến trúc tổng thể cho KMA Chatbot với hai phân hệ chính:
    \begin{itemize}
        \item \textit{Phân hệ Hỏi đáp (Chatbot):} Có khả năng giao tiếp tự nhiên với người dùng cuối, giải đáp chính xác các câu hỏi phổ biến về chỉ tiêu, điểm chuẩn, học phí, phương thức xét tuyển. Đặc biệt, hệ thống có khả năng tự động truy xuất tài liệu (Quy chế, biểu mẫu) thông qua RAG để sinh ra câu trả lời có tính cập nhật cao.
        \item \textit{Phân hệ Quản trị (Admin Dashboard):} Cung cấp giao diện trực quan cho người quản trị để quản lý tài liệu ngữ cảnh (Context Documents), theo dõi thống kê số lượng hội thoại, quản lý người dùng và phân quyền (RBAC) một cách hiệu quả.
    \end{itemize}
    
    \item \textbf{Về mặt ứng dụng thực tiễn:} Hệ thống đã chứng minh được tính khả thi trong việc giảm tải áp lực cho đội ngũ tư vấn viên trong mùa tuyển sinh. Cơ sở dữ liệu và kho tài liệu ngữ cảnh được tổ chức đa tầng (Qdrant, Redis) giúp tối ưu hóa thời gian phản hồi, mang lại trải nghiệm mượt mà cho thí sinh và phụ huynh.
\end{itemize}

\section{Hạn chế của hệ thống}
Mặc dù đã đạt được những kết quả nhất định, hệ thống KMA Chatbot vẫn còn tồn tại một số hạn chế cần được khắc phục trong tương lai:

\begin{itemize}
    \item \textbf{Khả năng hiểu ngữ cảnh phức tạp:} Dù đã tích hợp RAG, hệ thống đôi khi vẫn gặp khó khăn trong việc xử lý các câu hỏi mang tính suy luận logic phức tạp hoặc các câu hỏi kết hợp nhiều ý định cùng lúc (Multi-intent). LLM đôi khi vẫn sinh ra câu trả lời thiếu tự nhiên nếu ngữ cảnh truy hồi (retrieved context) từ Vector Database không đủ độ chính xác.
    \item \textbf{Dữ liệu huấn luyện chưa phong phú:} Tập dữ liệu huấn luyện cho Rasa NLU chủ yếu được xây dựng dựa trên kịch bản dự kiến và một lượng nhỏ dữ liệu thu thập thực tế. Để mô hình nhận diện tốt hơn các từ lóng, viết tắt hay sai chính tả của gen Z, hệ thống cần một tập dữ liệu lớn và đa dạng hơn nữa.
    \item \textbf{Phụ thuộc vào chất lượng tài liệu đầu vào:} Hệ thống RAG phụ thuộc rất lớn vào chất lượng tài liệu được đưa vào (Ingestion). Nếu tài liệu có định dạng bảng biểu phức tạp (như file Excel điểm chuẩn nhiều lớp) hoặc PDF scan không rõ nét, quá trình chia nhỏ (chunking) và trích xuất thông tin có thể bị sai lệch.
    \item \textbf{Hạn chế về hạ tầng triển khai:} Việc vận hành các Mô hình ngôn ngữ lớn (LLM) và Vector Database yêu cầu tài nguyên phần cứng (GPU/RAM) tương đối lớn, điều này tạo ra rào cản trong việc duy trì hiệu suất cao khi số lượng người dùng truy cập đồng thời tăng đột biến.
\end{itemize}

\section{Hướng phát triển tương lai}
Dựa trên những hạn chế đã phân tích, nhóm nghiên cứu đề xuất một số hướng phát triển tiếp theo để hoàn thiện và mở rộng hệ thống KMA Chatbot:

\begin{itemize}
    \item \textbf{Cải thiện chất lượng RAG và NLU:} 
    \begin{itemize}
        \item Ứng dụng các kỹ thuật tiên tiến hơn trong RAG như \textit{Graph RAG} (sử dụng cơ sở dữ liệu đồ thị tri thức - Knowledge Graph) để hệ thống có khả năng trả lời các câu hỏi mang tính suy luận, liên kết nhiều tài liệu khác nhau.
        \item Tích hợp cơ chế \textit{Re-ranking} trong quá trình truy hồi vector để đảm bảo các đoạn văn bản liên quan nhất luôn được ưu tiên đưa vào LLM.
        \item Bổ sung dữ liệu hội thoại thực tế từ các kênh tư vấn của Học viện (Fanpage, Hotline) để liên tục tái huấn luyện (re-train) mô hình Rasa NLU.
    \end{itemize}
    
    \item \textbf{Mở rộng phạm vi nghiệp vụ:} 
    \begin{itemize}
        \item Mở rộng chatbot không chỉ hỗ trợ tư vấn tuyển sinh mà còn hỗ trợ sinh viên đang học tập tại trường (tư vấn học vụ, lịch thi, đăng ký tín chỉ, thủ tục hành chính).
        \item Tích hợp hệ thống chatbot lên các nền tảng phổ biến khác như Facebook Messenger, Zalo Official Account hoặc Telegram để người dùng dễ dàng tiếp cận hơn.
    \end{itemize}
    
    \item \textbf{Tối ưu hóa quản trị và vận hành:}
    \begin{itemize}
        \item Xây dựng cơ chế đánh giá tự động chất lượng câu trả lời (Automated Evaluation) để phát hiện sớm các trường hợp chatbot trả lời sai (hallucination).
        \item Phát triển tính năng phân tích cảm xúc (Sentiment Analysis) để chatbot có thể nhận biết thái độ của người dùng, từ đó đưa ra phản hồi mềm mỏng hơn hoặc tự động chuyển hướng cuộc gọi đến tư vấn viên con người (Human handoff) khi người dùng có dấu hiệu bức xúc hoặc yêu cầu tư vấn chuyên sâu.
    \end{itemize}
\end{itemize}

Tóm lại, KMA Chatbot là một bước đệm quan trọng trong quá trình chuyển đổi số công tác tư vấn tuyển sinh của Học viện Kỹ thuật Mật mã. Với những hướng phát triển rõ ràng, hệ thống hứa hẹn sẽ trở thành một trợ lý ảo toàn diện, thông minh và thân thiện trong tương lai.
